\chapter{Project 4: Pocket DJ}

\section{Overview}
In this project, you will turn your microcontroller into a handheld musical instrument: the \textbf{Pocket DJ}!
You will wire three pushbuttons, a mini speaker, and a 3-pixel RGB LED strip to your breadboard.
Pressing each button triggers a distinct musical note on the speaker and lights up a matching colored pixel on the LED strip.

Over the course of this project, you will:
\begin{itemize}
    \item Connect multiple tactile pushbuttons and read them using MicroPython
    \item Use PWM (Pulse Width Modulation) to generate different pitch frequencies on a speaker
    \item Control individual addressable WS2812B RGB LEDs to create synchronized visual pulses
    \item Program rhythm loops and sound sequences using Python lists and tuples
\end{itemize}

\begin{figure}[H]
\centering
    \maybeincludegraphics[width=.6\linewidth]{project_4/success.jpg}
    \caption{The assembled Pocket DJ ready to drop some beats!}
\end{figure}

\begin{tcolorbox}[colback=blue!5!white,colframe=blue!40!black]
    \textbf{Learning goals:}
    \begin{itemize}
        \item Read pushbuttons with active-low logic and internal pull-up resistors
        \item Generate musical notes by controlling PWM frequency on a speaker
        \item Control WS2812B addressable LEDs using the \texttt{neopixel} library
        \item Structure musical sequences and automated rhythm loops in code
    \end{itemize}
\end{tcolorbox}

\pagebreak

\section{Components}
\begin{components}
    \item Seeed XIAO ESP32-C3 microcontroller
    \item 3-pixel WS2812B LED strip
    \item 3x Tactile pushbuttons ($6\times6$\,mm)
    \item 1x Treedix 1W 8\si{\ohm} mini speaker (with breadboard-friendly Dupont leads)
    \item Breadboard
    \item Jumper wires
\end{components}

\section{Directions}

\subsubsection{Remove previous components}
Before beginning, remove components and wires from the previous project. You may leave the
microcontroller seated in the breadboard.

\subsection{Creating the circuit}

% Seating the microcontroller. This is identical for every project, so the
% coordinates live here rather than in any one project chapter. The diagram
% generator reads the \bbhole definitions below for every breadboard diagram
% it draws, which is why this file is the only place they appear.
\subsubsection{Attach the microcontroller to the breadboard}
Orient the microcontroller so that the pin labeled \textbf{5V} goes into hole \bbhole{mcu.5v}{H1}
(Column H, Row 1), \textbf{GPIO2} into \bbhole{mcu.gpio2}{D1}, \textbf{GPIO20} into
\bbhole{mcu.gpio20}{H7}, and \textbf{GPIO21} into \bbhole{mcu.gpio21}{D7}. Then press it down---this
takes more force than you might expect.

\begin{figure}[H]
    \centering
    \maybeincludegraphics[width=.95\linewidth]{common/microcontroller_seated_in_breadboard.png}
    \caption{So far, so good!}
\end{figure}


\subsubsection{Ground rail distribution}
Because we have multiple components requiring a ground connection, we connect the microcontroller's
\textbf{GND} pin to the breadboard's top blue ($-$) power rail. This powers the entire rail with ground:

\begin{center}
\begin{tabular}{|l|l|l|}
    \hline
    \textbf{Connection} & \textbf{XIAO pin} & \textbf{Jumper holes} \\
    \hline
    \connectionrow{mcu-ground}{Top GND rail}{GND}{I2}{top-2}
    \hline
\end{tabular}
\end{center}

\subsubsection{WS2812B LED strip placement and wiring}
The LED strip has a right-angle header so it stands upright once seated.
Insert those three pins into row \textbf{F} (rows 22--24) on the right side of the breadboard so it is
clearly visible and not obscured by jumper wires:
\textbf{5V} into \bbhole{pixels.header.power}{F22},
\textbf{DI} into \bbhole{pixels.header.data}{F23}, and
\textbf{GND} into \bbhole{pixels.header.ground}{F24}.
Then connect the jumper wires:

\begin{center}
\begin{tabular}{|l|l|l|}
    \hline
    \textbf{WS2812B pin} & \textbf{Source / Pin} & \textbf{Jumper holes} \\
    \hline
    \connectionrow{pixels-power}{5V}{3V3}{I3}{J22}
    \connectionrow{pixels-data}{DI}{GPIO 20}{I7}{J23}
    WS2812B ground & Top GND rail & \bbhole{jumper.pixels-ground.start}{top-24} $\rightarrow$ \bbhole{jumper.pixels-ground.end}{J24} \\
    \hline
\end{tabular}
\end{center}

\subsubsection{Tactile button placement and wiring}
Seat the three pushbuttons in the lower half of the breadboard spanning columns \textbf{A} through \textbf{D}:
\begin{itemize}
    \item \textbf{Pad 0 (Loop/Bass):} Place one pair of legs in \bbhole{pad0.side-a}{A10} and \textbf{D10}, and the other pair in \bbhole{pad0.side-b}{A12} and \textbf{D12}.
    \item \textbf{Pad 1 (Mid note):} Place one pair of legs in \bbhole{pad1.side-a}{A14} and \textbf{D14}, and the other pair in \bbhole{pad1.side-b}{A16} and \textbf{D16}.
    \item \textbf{Pad 2 (High note):} Place one pair of legs in \bbhole{pad2.side-a}{A18} and \textbf{D18}, and the other pair in \bbhole{pad2.side-b}{A20} and \textbf{D20}.
\end{itemize}

Connect each button's signal leg (column \textbf{E}) to the microcontroller:
\begin{center}
\begin{tabular}{|l|l|l|}
    \hline
    \textbf{Button signal} & \textbf{XIAO pin} & \textbf{Jumper holes} \\
    \hline
    \connectionrow{pad0-signal}{Pad 0 (row 10)}{GPIO 21}{C7}{E10}
    \connectionrow{pad1-signal}{Pad 1 (row 14)}{GPIO 7}{C6}{E14}
    \connectionrow{pad2-signal}{Pad 2 (row 18)}{GPIO 6}{C5}{E18}
    \hline
\end{tabular}
\end{center}

Next, wire each button's ground leg (column \textbf{E}) to the top ground rail:
\begin{center}
\begin{tabular}{|l|l|l|}
    \hline
    \textbf{Button ground} & \textbf{Source} & \textbf{Jumper holes} \\
    \hline
    Pad 0 ground & Top GND rail & \bbhole{jumper.pad0-ground.start}{top-12} $\rightarrow$ \bbhole{jumper.pad0-ground.end}{E12} \\
    Pad 1 ground & Top GND rail & \bbhole{jumper.pad1-ground.start}{top-16} $\rightarrow$ \bbhole{jumper.pad1-ground.end}{E16} \\
    Pad 2 ground & Top GND rail & \bbhole{jumper.pad2-ground.start}{top-20} $\rightarrow$ \bbhole{jumper.pad2-ground.end}{E20} \\
    \hline
\end{tabular}
\end{center}

\subsubsection{Speaker placement and wiring}
Insert the speaker's two Dupont leads into the lower half of the breadboard:
the positive red lead into \bbhole{speaker.lead.signal}{C26} and the negative black lead into \bbhole{speaker.lead.ground}{C28}.
Then connect the speaker signal and ground jumpers:

\begin{center}
\begin{tabular}{|l|l|l|}
    \hline
    \textbf{Speaker lead} & \textbf{Source / Pin} & \textbf{Jumper holes} \\
    \hline
    \connectionrow{speaker-signal}{Signal (+)}{GPIO 10}{I4}{E26}
    Speaker ground ($-$) & Top GND rail & \bbhole{jumper.speaker-ground.start}{top-28} $\rightarrow$ \bbhole{jumper.speaker-ground.end}{E28} \\
    \hline
\end{tabular}
\end{center}

\begin{figure}[H]
    \centering
    \maybeincludegraphics[width=.95\linewidth]{project_4/wiring.png}
    \caption{Complete wiring diagram for Project 4: Pocket DJ.}
\end{figure}

\subsection{Programming the microcontroller}
Open the file named \texttt{project\_4\_pocket\_dj.py} in the IDE and run it.

When the script starts, the serial console displays:
\begin{verbatim}
Project 4: Pocket DJ
Pad 0: GPIO 21 | Pad 1: GPIO 7 | Pad 2: GPIO 6
Speaker: GPIO 10 | LED Strip: GPIO 20
Press any button to play music!
\end{verbatim}

Try pressing each pad:
\begin{itemize}
    \item \textbf{Pad 1} plays a middle note (330\,Hz) and flashes the middle LED green.
    \item \textbf{Pad 2} plays a higher note (392\,Hz) and flashes the third LED blue.
    \item \textbf{Pad 0} plays an automated 4-step rhythm loop across the notes and colors!
\end{itemize}

\subsection{How the program works}

\subsubsection{Musical tones with PWM}
In Project 1, PWM was used to control LED brightness by changing duty cycle at a fixed frequency.
Here, we use PWM on the speaker pin, but instead of changing brightness, we change the \textbf{frequency} (the pitch):
\begin{lstlisting}[language=Python, caption=Playing a tone on the speaker]
pwm = PWM(Pin(PIN_SPEAKER), freq=freq, duty=512)
time.sleep_ms(duration_ms)
pwm.deinit()
\end{lstlisting}
A higher frequency makes the speaker vibrate faster, producing a higher musical pitch. A 50\% duty cycle (\texttt{duty=512}) provides optimal volume. Calling \texttt{pwm.deinit()} silences the speaker when the note is done.

\subsubsection{RGB LEDs with NeoPixel}
The WS2812B LEDs are addressable: a single digital data pin (GPIO 20) controls the color and brightness of all three pixels independently.
Colors are specified as \texttt{(Red, Green, Blue)} tuples from 0 to 255:
\begin{lstlisting}[language=Python, caption=Lighting a pixel]
pixels[pixel_idx] = color
pixels.write()
\end{lstlisting}

\subsubsection{Rhythm loops}
Pad 0 reads a sequence of note indices and timing delays stored in a Python tuple:
\begin{lstlisting}[language=Python, caption=Rhythm loop sequence]
LOOP = (
    (0, 80),   # Pad 0, rest 80 ms
    (2, 60),   # Pad 2, rest 60 ms
    (1, 80),   # Pad 1, rest 80 ms
    (2, 120),  # Pad 2, rest 120 ms
)
\end{lstlisting}
When triggered, a \texttt{for} loop steps through each element, playing each note and waiting before the next beat.

\subsection{Hardware test checklist}
\begin{enumerate}
    \item Press Pad 1: You should hear a tone and see the green LED light up.
    \item Press Pad 2: You should hear a higher tone and see the blue LED light up.
    \item Press Pad 0: You should hear the 4-step beat loop play across the different tones and LEDs.
    \item Check serial output in the IDE to verify that button presses are detected cleanly without double triggers.
\end{enumerate}

\begin{tcolorbox}[colback=yellow!10!white,colframe=yellow!50!black]
    \textbf{If something does not work:}
    \begin{itemize}
        \item \textbf{No sound from speaker:} Check that the speaker red lead is in \bbref{speaker.lead.signal} and jumpered from \bbref{jumper.speaker-signal.start} to \bbref{jumper.speaker-signal.end} (GPIO 10). Check that the black lead is in \bbref{speaker.lead.ground} and grounded to the top rail via \bbref{jumper.speaker-ground.start} to \bbref{jumper.speaker-ground.end}.
        \item \textbf{LED strip does not light:} Check that the strip is seated in F22--F24 with DI in \bbref{pixels.header.data} and connected to GPIO 20 (\bbref{jumper.pixels-data.start} to \bbref{jumper.pixels-data.end}), power to 3.3V (\bbref{jumper.pixels-power.start} to \bbref{jumper.pixels-power.end}), and GND to the top rail (\bbref{jumper.pixels-ground.start} to \bbref{jumper.pixels-ground.end}).
        \item \textbf{Button press does nothing:} Verify button orientation and connections. Check that the microcontroller GND wire (\bbref{jumper.mcu-ground.start} to \bbref{jumper.mcu-ground.end}) is connected to the top rail, and each button ground jumper is connected.
    \end{itemize}
\end{tcolorbox}

\section{Review}
In this project, you constructed a multi-component digital instrument combining inputs (pushbuttons) with two different types of outputs (PWM audio on a speaker and addressable RGB LEDs). You learned how frequency determines pitch, how to drive NeoPixel LEDs, and how to program sequenced loops.

\section{Possible Extensions}
\begin{itemize}
    \item \textbf{Custom melodies:} Modify the \texttt{LOOP} tuple in \texttt{project\_4\_pocket\_dj.py} to create your own custom beats or jingles.
    \item \textbf{Scale selector:} Add higher or lower octaves by doubling or halving the frequencies in \texttt{PADS}.
    \item \textbf{DJ light show:} Update \texttt{play\_note} to create visual ripple effects where non-playing LEDs pulse with complementary colors.
\end{itemize}
